\reportheading{2. Cooling load and airflow distribution (Task 1)}
\subhead{2.1 Preliminary thermal balance}
The selected base design uses an office cooling allowance of 189 W/m² and an approximate sensible heat factor of 0.90. Since BEC 2024 [6, p. 85] uses an 11 K room-to-supply temperature difference for calculating the design supply airflow of the reference building, this difference is applied to the office to select a design supply air temperature of $24-11=13$ °C. Table 2 uses this selected value to calculate the sensible heat balance.

\tabletitle{Table 2. Thermal balance and supply airflow.}
\begingroup\setlength{\tabcolsep}{3bp}
\begin{reporttable}{L{0.629630}L{0.370370}}
\rowcolor{tablegrey} \textbf{Calculation} & \textbf{Result} \\\hline
Total cooling = 1022 x 189 & 193.158 kW thermal \\
Sensible cooling = 193.158 x 0.90 & 173.8422 kW thermal \\
$Q=\frac{q_s}{\rho c_p(24-13)}$ & 13091.316 L/s \\
Minimum supply = 0.40 x maximum & 5236.526 L/s \\
\end{reporttable}\endgroup
At the selected 11 K temperature difference, each cubic metre per second of supply air removes about 13.28 kW of sensible heat. Therefore, the office needs 13.09 m³/s to carry its 173.84 kW sensible load. The remaining cooling duty is latent and is not included in this airflow calculation.

\subhead{2.2 Diffuser distribution and VAV grouping}
Larger service areas are assigned a larger share of floor airflow, and the four outlet flows on each box add up to its design airflow (Appendix C).

Each of the 80 diffusers fits inside a distinct ceiling tile within its assigned area, as shown in Figure 3.


The service-cell areas are scaled to the annotated floor area before allocating the sensible cooling load. For SD-77, the allocated area is 15.7843 m², and its peak flow is $\frac{13091.316\times15.7843}{1022}=202.19$ L/s, with Appendix A listing the resulting different outlet flows.

Dividing the floor total equally gives 163.64 L/s per outlet. While this average works for an estimate, it would not provide enough air for large zones such as SD-77. By scaling airflow by area, SD-77 requires about 38.55 L/s more than the average without changing the total floor airflow. As a result, this difference changes both the terminal-box airflow and the peak throw at this outlet.

\subhead{2.3 Characteristic length}
In each of the four discharge directions, the characteristic length is defined as the shorter distance to either the office boundary or the midpoint between opposing outlets, and these four lengths are listed in Appendix B. For the facing pair SD-77 and SD-78, the 3.000 m centre spacing gives L = 1.500 m in the opposing directions, while the same outlet has a separate length north, south and west. Changing the service area changes its airflow and throw, whereas changing the tile position changes these lengths even at the same flow. Appendix A and Appendix B record the two inputs separately.
